\section{Profilazione utente}
\paragraph{\textbf{Chi è l’utente?}}
Il rifugio opera in Italia ed è concepito per utenti stabiliti sul territorio nazionale (nota: il numero di cellulare non presenta restrizioni geografiche).  \\
Gli utenti non hanno limiti di età particolari: il sito web utilizza un linguaggio semplice, adatto anche ai più giovani, che si presume possano essere tra i principali motivatori nell’indurre i genitori ad adottare un animale.  \\
È stata quindi considerata una categoria di utenti che spazia da persone con buone competenze digitali a persone con capacità limitate nell’utilizzo delle tecnologie informatiche.

\paragraph{\textbf{Che bisogni ha?}}
L’utente accede al sito principalmente per visualizzare la disponibilità degli animali e per comprendere le modalità di adozione.  \\
Di conseguenza, queste informazioni sono messe in particolare evidenza all’interno del sito.  \\
Un ulteriore bisogno può essere quello di entrare in contatto con il rifugio; sono disponibili: un modulo di contatto, indirizzi email, numero di telefono e canali social.  \\
È stata inoltre considerata una categoria di utenti, seppur numericamente inferiore, interessata ad affiliarsi al rifugio, sia come professionisti che come volontari.

\paragraph{\textbf{Che competenze possiede?}}
Non è richiesta alcuna competenza specifica per navigare all’interno del sito, dato il bacino d’utenza molto generale a cui è rivolto.

\paragraph{\textbf{In che contesto lo usa?}}
Considerata la varietà dell’utenza, sono stati previsti diversi contesti di utilizzo.\\
Il sito può essere consultato sia da computer che da dispositivi mobili. 
È stata inoltre predisposta una modalità di stampa pensata per adattarsi alla consultazione su supporto cartaceo, risultando particolarmente utile per utenti che non utilizzano direttamente dispositivi digitali.

